\section{Explicit examples}
\label{sec:examples}

Three models illustrate the $\Ainf$ operation tower at successive depths: the free chiral multiplet (Example~\ref{rem:free_chiral}) with $m_{k \ge 3} = 0$; Landau--Ginzburg with cubic superpotential (Example~\ref{rem:LG}) with $m_3 \neq 0$; and abelian Chern--Simons (Example~\ref{rem:abelian_CS}) with boundary algebra $\widehat{\mathfrak{u}(1)}_k$.

In each model the BV data determine $Q=m_1$, the propagator determines
$m_2$, the Feynman graphs determine $m_{k\geq 3}$, and the cohomology
carries the Poisson vertex structure.

\begin{remark}[Analytic hypotheses]
The calculations use the hypotheses of Theorem~\ref{thm:physics-bridge}:
BV product gauge, a factorized logarithmic parametrix, one-loop finiteness,
and product-polynomial interactions. Under these hypotheses the resulting
observables form a logarithmic $\SCchtop$-algebra, and
Theorem~\ref{thm:FM-calculus} supplies the logarithmic-form, residue, and
boundary-factorization identities used below. The free multiplet and
abelian Chern--Simons examples are Gaussian, so only the parametrix enters
the explicit formulas.
\end{remark}

\subsection{Example I: free chiral multiplet}
\label{rem:free_chiral}

All $m_{k \geq 3}$ vanish identically.

\subsubsection{Field Content and Action}

A free chiral multiplet on $\R \times \C$ consists of:
\begin{itemize}
\item A complex scalar $\phi : \R \times \C \to \C$ (the bottom component, in cohomological degree~$0$);
\item A fermionic field $\psi$ (the BV antifield, in cohomological degree~$1$).
\end{itemize}

In the BV formalism after gauge fixing, the action is
\[
S_{\text{free}} = \int_{\R \times \C} \langle \psi, (d_t + \bar{\partial})\phi \rangle \, dt \, dz \, d\bar{z},
\]
where $d_t + \bar{\partial}$ is the HT kinetic operator.

\subsubsection{Local Observables}

The algebra of local observables is the polynomial algebra
\[
\A_{\text{free}} = \C[\phi_n, \psi_n \mid n \geq 0],
\]
generated by derivatives
\[
\phi_n := \frac{1}{n!} \partial_z^n \phi, \quad \psi_n := \frac{1}{n!} \partial_z^n \psi,
\]
with cohomological degrees $|\phi_n| = 0$, $|\psi_n| = 1$.

\subsubsection{BV-BRST Differential: Explicit Formulas}

\begin{proposition}[Free BV-BRST differential]
\label{prop:free_BRST_detailed}
The BV-BRST differential $Q = \{S_{\text{free}}, -\}_{\mathrm{BV}}$ acts on generators by
\begin{align}
Q \cdot \phi_n &= \psi_n, \label{eq:Qphi_free}\\
Q \cdot \psi_n &= 0. \label{eq:Qpsi_free}
\end{align}
It extends to all of $\A_{\text{free}}$ by the graded Leibniz rule $Q(a\cdot b) = (Qa)\cdot b + (-1)^{|a|}a\cdot (Qb)$.
\end{proposition}

\begin{proof}
The BV bracket pairs the field~$\phi$ with the antifield~$\psi$ via $\{\phi(x),\psi(y)\}_{\mathrm{BV}} = \delta^{(3)}(x-y)$. The functional derivative of $S_{\text{free}}$ with respect to the antifield~$\psi$ is
\[
Q\phi = \frac{\delta S_{\text{free}}}{\delta \psi} = (d_t + \bar{\partial}_z)\phi.
\]
Passing to holomorphic modes $\phi_n = \frac{1}{n!}\partial_z^n\phi$ gives $Q\phi_n = \psi_n$. Since $\psi$ is the BV antifield (of top cohomological degree in the two-term complex), $Q\psi_n = 0$. Then
\[
Q^2\phi_n = Q(\psi_n) = 0,
\]
so $(\A_{\text{free}}, Q)$ is a cochain complex. The Leibniz rule follows from the derivation property of the BV bracket.
\end{proof}

\begin{remark}[Free $\neq$ Trivial Differential]
The theory is \emph{free} (no interaction vertices), but the differential $Q$ is nonzero. It encodes the classical equations of motion via the BV antifield pairing. ``Free theory'' means $m_{k\ge 3}=0$, not $Q=0$.
\end{remark}

\subsubsection{The $\lambda$-Bracket $m_2$}

The binary operation $m_2$ is governed by the free propagator.

\begin{construction}[Free propagator on $\R \times \C$]
\label{const:free_propagator_detailed}
The Green's function for $(d_t + \bar{\partial}_z)$ is
\[
K(t-t',z-z') = \frac{1}{2\pi} \frac{\Theta(t-t')}{z-z'},
\]
with $\Theta$ the Heaviside step function. It produces the unique distributional solution of $(d_t+\bar{\partial}_z)K = \delta(t-t')\delta^{(2)}(z-z')$.
\end{construction}

\begin{proposition}[Explicit $\lambda$-brackets]
\label{prop:free_lambda_brackets_detailed}
The only nonvanishing $\lambda$-bracket is
\begin{align}
\label{eq:free_lambda_phi_psi}
\{\phi_n \,{}_\lambda\, \psi_m\} &= \binom{n+m}{n} (-1)^n \lambda^{n+m},
\end{align}
with all others obtained from sesquilinearity and the $Q$-derivation property.
\end{proposition}

\begin{proof}
Contraction with the propagator gives
\[
\{\phi(z_1)\,{}_\lambda\,\psi(z_2)\} = \frac{\Theta(t_1-t_2)}{2\pi(z_1-z_2)}.
\]
Expansion of $e^{\lambda(z_1-z_2)}$ and differentiation $n$ and $m$ times in $z_1,z_2$ produces \eqref{eq:free_lambda_phi_psi}. Sesquilinearity determines the remaining brackets.
\end{proof}

\begin{remark}[Sesquilinearity check]
Equation \eqref{eq:free_lambda_phi_psi} immediately implies
\[
\{\partial \phi_n \,{}_\lambda\, \psi_m\} = -\lambda \{\phi_n \,{}_\lambda\, \psi_m\},\qquad
\{\phi_n \,{}_\lambda\, \partial \psi_m\} = (\partial+\lambda)\{\phi_n \,{}_\lambda\, \psi_m\},
\]
verifying the PVA sesquilinearity relations.
\end{remark}

\subsubsection{$m_{k \geq 3}$ Vanish at Chain Level}

\begin{proposition}[No higher operations in free theories]
\label{prop:free_mk_zero}
For the free chiral multiplet,
\[
m_k = 0 \quad \text{for all } k \geq 3.
\]
\end{proposition}

\begin{proof}
Higher operations $m_k$ arise from Feynman diagrams with interaction vertices. Since $S_{\text{free}}$ contains no interaction terms (quadratic in fields), all Feynman diagrams contributing to $m_{k \geq 3}$ vanish identically: no vertices for internal points, no diagrams.
\end{proof}

\begin{remark}[Koszulness from clean self-intersection]
\label{rem:free-koszul-clean}
The free multiplet is Koszul because its BV action is
quadratic: all Feynman diagrams with interaction vertices
vanish, giving $m_k = 0$ for $k \ge 3$
(Proposition~\ref{prop:free_mk_zero}). Geometrically, the
Lagrangian self-intersection
$\cL \times_{\cM}^h \cL$ is clean (no higher Tor), so the
self-intersection groupoid is \'etale and bar-cobar inversion
holds without completion.
In the depth classification of
\S\ref{subsec:excess-intersection-depth}, this is
class~$\mathbf{G}$ (depth~$2$): the shadow obstruction tower terminates
at degree~$2$ because the abelian field content admits no
nontrivial higher graph sums.
\end{remark}

\subsubsection{Cohomology Computation}

\begin{proposition}[Cohomology of the free chiral multiplet]
\label{prop:free_cohomology}
We have $H^0(\A_{\text{free}}, Q) = \C$ and $H^k(\A_{\text{free}}, Q) = 0$ for $k \neq 0$.
\end{proposition}

\begin{proof}
Each pair $(\phi_n,\psi_n)$ forms an acyclic $Q$-doublet
\[
\phi_n \xrightarrow{Q} \psi_n \xrightarrow{Q} 0,
\]
so the nontrivial cohomology class is represented by the constant~$1$.
\end{proof}

\begin{remark}[Trivial Descended PVA]
Since $H^\bullet(\A_{\text{free}}, Q) = \C$, the induced Poisson vertex algebra on cohomology is one-dimensional (the trivial PVA).
\end{remark}

\subsection{Example II: Landau--Ginzburg model with cubic superpotential}
\label{rem:LG}

$m_3 \neq 0$.

\subsubsection{Setup: LG Theory with $W(\phi) = \lambda_3 \phi^3/3$}

Consider a chiral multiplet $(\phi, \psi)$ with superpotential
\[
W(\phi) = \frac{\lambda_3}{3} \phi^3.
\]
Set $\lambda_3 = 1$ (the coupling is restored by dimensional analysis $\phi \mapsto \lambda_3^{1/2}\phi$); this matches the normalisation $W = \phi^3/3$ of \S\ref{lem:LG_cubic_complete}.
The BV action is
\[
S_{\text{LG}} = S_{\text{free}} + \int_{\R \times \C} W'(\phi) \psi \, dt \, dz \, d\bar{z} = S_{\text{free}} + \int_{\R \times \C} \lambda_3 \phi^2 \psi \, dt \, dz \, d\bar{z}.
\]
The interaction term $\phi^2 \psi$ is the \emph{interaction vertex} generating higher operations.

\subsubsection{BV-BRST Differential}

The differential now includes interaction contributions from the superpotential:
\begin{align}
Q \cdot \phi &= \psi + W'(\phi) = \psi + \lambda_3 \phi^2, \\
Q \cdot \psi &= 0.
\end{align}
Nilpotency $Q^2 = 0$ follows from the classical master equation $\{S_{\text{LG}}, S_{\text{LG}}\}_{\mathrm{BV}} = 0$. Restricted to the generators alone the computation $Q^2\phi = Q(\psi + \lambda_3\phi^2) = 2\lambda_3\phi(\psi + \lambda_3\phi^2)$ does not vanish termwise; $Q^2 = 0$ holds on the full BV complex, where the antifield pairings cancel the residual terms. The structural argument is $Q = \{S_{\mathrm{LG}}, -\}_{\mathrm{BV}}$ with $Q^2 = \{\tfrac{1}{2}\{S,S\}, -\}_{\mathrm{BV}} = 0$ by the graded Jacobi identity for the antibracket.

\begin{remark}
More precisely, on the algebra of local observables
\(\A_{\text{LG}} = \C[\phi_n, \psi_n]\), the BV differential is
\(Q_{\BV}\phi_n = \psi_n\) on the free doublets, and the
superpotential deformation enters through the transferred higher
operations. After passing to the matrix-factorisation model one uses
the Koszul differential \(d_W=\partial W/\partial\phi\); in that model
\(m_1=d_W\) and the cohomology is the Jacobian ring.
\end{remark}

\subsubsection{The Higher Operation $m_3$}

\begin{proposition}[Cubic $m_3$ from the superpotential]
\label{prop:LG_m3}
For the LG model with cubic \(W\), the transferred ternary operation is
the polarized third Taylor coefficient of the superpotential. On the
Jacobian model,
\begin{equation}
\label{eq:LG_m3_formula}
m_3(O_1,O_2,O_3)=
W'''(\phi)\,\partial_\phi O_1\,\partial_\phi O_2\,\partial_\phi O_3
=2\lambda_3\,\partial_\phi O_1\,\partial_\phi O_2\,\partial_\phi O_3 .
\end{equation}
\end{proposition}

\begin{proof}[Computation]
The interaction vertex is the cubic Taylor coefficient of \(W\).
There is no internal edge in the primitive ternary tree, hence no
spectral denominator. The residue at the total-collision stratum is
the local coefficient \(W'''=2\lambda_3\). Rescaling the BV pairing
rescales the propagator and the vertex together.
\end{proof}

\subsubsection{$m_{k \geq 4}$}

\begin{proposition}[Truncation at $m_3$ for cubic Landau--Ginzburg]
\label{prop:LG_truncation}
On the Jacobian minimal model of the cubic Landau--Ginzburg theory,
the operations satisfy:
\begin{itemize}
\item \(m_1^{J}=0\);
\item \(m_2^{J}\) is the product on \(J(W)\);
\item \(m_3^{J}\) is the cubic Taylor coefficient~\eqref{eq:LG_m3_formula};
\item \(m_k^{J}=0\) for all \(k\geq4\).
\end{itemize}
\end{proposition}

\begin{proof}
The transferred minimal model is controlled by the Taylor coefficients
of \(W\).

\textbf{Cubic vertex constraint.} The primitive \(r\)-ary operation is
the polarization of the \(r\)-th Taylor coefficient \(W^{(r)}\), followed
by projection to the Jacobian ring \(J(W)\). For
\(W=\lambda_3\phi^3/3\),
\[
W'''=2\lambda_3,\qquad W^{(r)}=0\quad(r\geq4).
\]
A tree with several cubic vertices contributes to a Stasheff composite
of \(m_2\) and \(m_3\), not to a new primitive \(m_k\). Loop graphs have
positive BV degree and do not survive in the genus-zero minimal
operation on \(J(W)\).

The nonvanishing operations are:
\begin{itemize}
\item \(m_1^{J}=0\);
\item \(m_2^{J}\), the induced multiplication;
\item \(m_3^{J}\), the single cubic Taylor vertex.
\end{itemize}
Thus \(m_k^{J}=0\) for \(k\geq4\).
\end{proof}

\begin{remark}[General Polynomial Superpotentials]
For a general polynomial superpotential $W(\phi) = \sum_{d=2}^D \lambda_d \phi^d/d$, the highest-degree term $\lambda_D\phi^D$ supplies a $D$-valent vertex $W^{(D)} = (D-1)!\,\lambda_D$, so the operation $m_D$ on the Jacobian ring $J(W) = \C[\phi]/(W')$ is nonzero. The cubic truncation \(m_{k}^{J}=0\) for \(k\geq4\) of Proposition~\ref{prop:LG_truncation} is special to \(D=3\): there \(W'''=2\lambda_3\) is constant and \(W^{(k)}=0\) for \(k\geq4\), so no primitive \(k\)-ary Taylor vertex exists. For $D \ge 4$ the operations $m_k$ persist to order~$D$ in the minimal $A_\infty$ model of $J(W)$.
\end{remark}

\subsubsection{Evaluation of $m_3$}

The ternary operation $m_3$ is constant in the spectral parameters and equals the third derivative of the superpotential.

\begin{computation}[Explicit $m_3$ for Landau--Ginzburg]
\label{comp:lg-m3-explicit}
For the LG model with superpotential \(W(\phi)=\lambda_3\phi^3/3\) on
\(\C\times\R\), the minimal \(A_\infty\) structure on \(J(W)\) is:
\begin{enumerate}[label=(\roman*)]
\item \emph{Differential.}
 \(m_1^{J}=0\). The differential before transfer is the Koszul
 differential \(d_W=\partial W/\partial\phi=\lambda_3\phi^2\).

\item \emph{Binary product.}
 \(m_2^{J}(a,b)=a\cdot b\) (the free commutative product; no $\lambda$-bracket
 corrections at tree level).

\item \emph{Ternary operation.}
 \begin{equation}\label{eq:lg-m3-evaluated}
 m_3^{J}(a,b,c;\lambda_1,\lambda_2)
 \;=\; W'''(\phi)\,\partial_\phi a\,\partial_\phi b\,\partial_\phi c
 \;=\; 2\lambda_3\,\partial_\phi a\,\partial_\phi b\,\partial_\phi c.
 \end{equation}
 The result is \emph{constant} in the spectral parameters $\lambda_1,\lambda_2$:
 the sole cubic vertex $W'''=2\lambda_3$ contributes a local interaction with no
 internal propagators, hence no spectral-parameter dependence.

\item \emph{Higher operations.}
 \(m_k^{J}=0\) for all \(k\geq4\), since \(W^{(k)}=0\) for \(k\geq4\)
 (Proposition~\ref{prop:LG_truncation}).
\end{enumerate}
\end{computation}

\begin{proof}
The Feynman diagram contributing to $m_3$ is the unique tree with a single cubic vertex ($V=1$, $E=0$ internal edges) and three external legs. Its vertex factor is $W'''(\phi)=2\lambda_3$.

\textbf{Cubic tree.}
A connected tree with $k$ external legs and cubic vertices has $V = k-2$, $E = k-3$. For $k=3$: $V=1$, $E=0$; no internal propagators, so the Feynman integral collapses to the vertex factor.

\textbf{Spectral independence.}
With no internal edges, no propagator denominator enters the primitive
operation, and no spectral-parameter dependence is generated. The
total-collision residue is the constant \(W'''=2\lambda_3\).

\textbf{Stasheff relation.}
The \(n=3\) Stasheff relation is associativity of \(m_2^{J}\).
The first relation containing \(m_3^{J}\) is the Hochschild cocycle
condition. The polarized cubic cochain is normalized; on the reduced
basis it is nonzero only on \((\phi,\phi,\phi)\), the middle terms
contain \(\phi^2\), and the two boundary multiplications cancel with
the suspended signs.
\end{proof}

\begin{remark}[Physical meaning of $m_3$]
\label{rem:lg-m3-physics}
The operation \(m_3=W'''\,\partial a\,\partial b\,\partial c\) is the
cubic interaction vertex of the Landau--Ginzburg model. Higher trees
are Stasheff composites of
\(m_2\) and \(m_3\); the minimal model has no independent \(k\)-ary
operation for \(k\geq4\).
\end{remark}

\begin{remark}[Connection to the Rosetta Stone]
\label{rem:lg-m3-rosetta}
The Swiss-cheese structure (\S\ref{subsec:rosetta-operations}) produces higher operations beyond the formal (Heisenberg) case, where all $m_{k \ge 3}$ vanish by formality (Remark~\ref{rem:rosetta-formality-koszulness}). The ternary operation $m_3 = W'''$ is the closed-color (holomorphic) contribution from the cubic interaction; it is constant in the spectral parameters $\lambda_1,\lambda_2$ since the vertex has no internal propagators. In the language of Section~\ref{subsec:rosetta-operations}, the bar element $\bar{m}_3 \in \bar{B}^3(\A_{\mathrm{LG}})$ is a nonzero cocycle representing the \emph{essential} non-formality of the LG model.
\end{remark}

\subsubsection{Cohomology and Descent to PVA}

\begin{proposition}[Landau--Ginzburg cohomology and Jacobian ring]
\label{prop:LG_cohomology}
The Jacobian ring of the superpotential is
\[
J(W) \;=\; \C[\phi]/(\partial W/\partial\phi)
 \;=\; \C[\phi]/(\lambda_3\phi^2),
\]
a $2$-dimensional algebra with basis $\{1,\phi\}$.
This is the cohomology $H^\bullet(\A_{\mathrm{LG}}, d_W)$: each pair
$(\phi_n,\psi_n)$ with $n \ge 2$ forms an acyclic $m_1$-doublet, while the
classes $1$ and $\phi$ survive.

The induced $A_\infty$ structure on $J(W)$ inherits a nonvanishing ternary
operation:
\begin{equation}\label{eq:lg-jacobian-m3}
 m_3^{J(W)}(\phi,\phi,\phi) = 2\lambda_3,
\end{equation}
the residual cubic coupling. The operation is strictly unital, hence it
vanishes as soon as one input is \(1\).
\end{proposition}

\begin{proof}
The Koszul differential \(d_W\) acts by multiplication by
\(\lambda_3\phi^2\). In the polynomial ring \(\C[\phi_n,\psi_n]\),
the kernel on degree-$0$ generators modulo the image is
\(\C[\phi]/(\lambda_3\phi^2)\). The higher modes \(\phi_n\)
(\(n\geq2\)) pair with \(\psi_n\) in acyclic doublets, so
\[
H^\bullet(\A_{\mathrm{LG}},d_W)\cong J(W).
\]

The homotopy transfer theorem gives
\(m_3^{J(W)} = \pi \circ m_3 \circ
(\iota \otimes \iota \otimes \iota)\), where
\(\iota\colon J(W) \hookrightarrow \A_{\mathrm{LG}}\) is the inclusion
and \(\pi\colon \A_{\mathrm{LG}} \twoheadrightarrow J(W)\) the
projection. By Computation~\ref{comp:lg-m3-explicit}, the polarized
vertex consumes one \(\phi\)-variation from each input, so
\(m_3^{J(W)}(\phi,\phi,\phi)=\pi(2\lambda_3)=2\lambda_3\). If an input
is \(1\), one derivative is zero.
\end{proof}

\begin{remark}[Chain-level vs.\ cohomology-level $m_3$]
The operation \(m_3^{J(W)}\) belongs to the matrix-factorisation
model with differential \(d_W\), not to the free \(Q_{\BV}\)-cohomology.
Thus \(m_3^{J(W)}(\phi,\phi,\phi)=2\lambda_3\) is the residue of the
superpotential, while the \(Q_{\BV}\)-cohomology sees only the free
doublet contraction.
\end{remark}

\subsection{Example III: abelian gauge theory with Chern--Simons coupling}
\label{rem:abelian_CS}

Gauge theories with Chern--Simons terms produce boundary chiral algebras.

\subsubsection{3d Abelian Gauge Theory on $\R \times \C$}

Consider $U(1)$ gauge theory on $\R \times \C$ with Chern--Simons coupling:
\[
S_{CS} = \frac{k}{4\pi} \int_{\R \times \C} A \wedge dA,
\]
where $A$ is a $U(1)$ gauge connection and $k \in \Z$ is the \emph{Chern--Simons level}.

After HT twist and BV gauge fixing, the theory has:
\begin{itemize}
\item Gauge field $A = A_z dz + A_{\bar{z}} d\bar{z} + A_t dt$;
\item Ghost fields $c$ (ghost), $\bar{c}$ (antighost);
\item Auxiliary fields completing the BV multiplet.
\end{itemize}

\subsubsection{Boundary conditions and chiral algebras}

Impose Dirichlet boundary conditions at $t = 0$, creating a boundary $\{t=0\} \cong \C$.

\begin{theorem}[Boundary chiral algebra for abelian Chern--Simons]
\label{thm:CS_boundary}
For $U(1)$ Chern--Simons theory at integral level $k$ on
$\R_{\geq 0}\times\C$ with Dirichlet boundary condition
$A_{\bar z}|_{t=0}=0$, write the flat boundary field as
$A_z=\partial\varphi$ and set \(J(z)=k\,\partial\varphi(z)\). Then
\(J\) generates
the rank-one Heisenberg vertex algebra $V^k(\mathfrak u(1))$.
Equivalently,
\[
J(z)J(w)\sim \frac{k}{(z-w)^2}.
\]
\end{theorem}

\begin{proof}
\textbf{Boundary field.}
Dirichlet boundary conditions $A_{\bar z}|_{t=0} = 0$ on the boundary
$\{t=0\} \cong \C$ leave the holomorphic component $A_z(0,z)$ as the
unconstrained boundary degree of freedom.  The flatness equation gives
$A_z=\partial\varphi$ on the boundary, and the Chern--Simons boundary
symplectic form is the standard chiral-boson form at level \(k\).
Thus
\[
\partial\varphi(z)\partial\varphi(w)\sim
\frac{1}{k(z-w)^2},
\qquad
J(z):=k\,\partial\varphi(z).
\]

\textbf{Affine cocycle.}
The resulting singular product is
\[
J(z)J(w)\sim \frac{k}{(z-w)^2}.
\]
Equivalently,
\begin{equation}
\label{eq:U1_OPE}
\{J_\lambda J\}=k\lambda.
\end{equation}
The bulk propagator
\[
\langle A_z(t_1,z_1)\, A_{\bar z}(t_2,z_2) \rangle
= \frac{k}{2\pi}\frac{\Theta(t_1 - t_2)}{z_1 - z_2}
\]
is the bulk kernel whose boundary symplectic residue gives this affine
two-cocycle; the double pole belongs to the current
\(J=k\partial\varphi\), not to the undifferentiated boundary potential.

\textbf{Abelian truncation.}
Abelian CS is a free (Gaussian) theory after gauge-fixing; connected Feynman diagrams with $\ge 3$ external legs vanish (no interaction vertices). Thus $m_k = 0$ for $k \ge 3$, and the boundary algebra is a strict vertex algebra.

\textbf{Non-abelian shadow.}
For non-abelian gauge group $G$ with Lie algebra $\mathfrak{g}$, the
moment-map current \(J^a\) generates $\widehat{\mathfrak{g}}_k$ with
OPE $J^a(z) J^b(w) \sim k\delta^{ab}/(z-w)^2 +
f^{ab}_c J^c(w)/(z-w)$, where $f^{ab}_c$ are structure constants.
The cubic vertex $\int f^{ab}_c A^a A^b A^c$ generates non-trivial
$m_k$ for $k \ge 3$; see \cite{CDG20,KZ25}.
\end{proof}

\subsubsection{Bulk PVA and bulk--boundary correspondence}

\begin{proposition}[Bulk PVA for abelian Chern--Simons under the physics bridge]
Assume the field-theoretic realization statement of
Theorem~\ref{thm:physics-bridge}. For the abelian HT Chern--Simons
realization, the bulk cohomology $H^\bullet(\A_{\mathrm{bulk}}, Q)$ is
a commutative Poisson vertex algebra, with commutativity forced by
topological invariance in~$t$ and with $\lambda$-bracket equal to the
classical Poisson bracket of the Chern--Simons action.
\end{proposition}

Under the same field-theoretic realization and exact-sector scope,
Theorem~\ref{thm:bulk_hochschild} identifies the abelian bulk
observable complex with the chiral Hochschild cochains of
$\widehat{\mathfrak{u}(1)}_k$, realising the framework of
Sections~\ref{sec:chiral_hochschild_expanded}--\ref{sec:bar_cobar}.

\begin{remark}[Connection to quantum groups]
The boundary $\widehat{\mathfrak{g}}_k$ is the classical shadow of $U_q(\widehat{\mathfrak{g}})$ at $q = e^{2\pi i/(k + h^\vee)}$. Line operators in the 3d theory are representations of $U_q(\widehat{\mathfrak{g}})$: the quantum group has a field-theoretic origin.
\end{remark}

\subsection{Summary Table}

\begin{table}[h]
\centering
\begin{tabular}{|l|c|c|c|}
\hline
\textbf{Property} & \textbf{Free Chiral} & \textbf{LG Cubic} & \textbf{Abelian CS} \\
\hline
Fields & $(\phi, \psi)$ & $(\phi, \psi)$ & $(A_z, c, \bar{c})$ \\
\hline
Interaction & None & $\lambda_3 \phi^3/3$ & $k\,A\wedge dA$ \\
\hline
$Q_{\BV}$ & $Q\phi=\psi$, $Q\psi=0$ & $Q_{\BV}\phi=\psi$ & Gauge BRST \\
\hline
Matrix-factorisation differential & -- & $d_W=\lambda_3\phi^2$ & -- \\
\hline
$m_2$ & Free propagator & Free propagator & Gauge propagator \\
\hline
$m_3$ & $0$ & $W'''=2\lambda_3$ & Gauge vertex \\
\hline
$m_{k\ge 4}$ & $0$ & $0$ (\(W^{(k)}=0\)) & Higher gauge ops \\
\hline
$H^\bullet$ & $\C$ & $H^\bullet(\A,d_W)=J(W)$ & Commutative PVA \\
\hline
Boundary algebra & -- & Matrix factorizations & $\widehat{\mathfrak{u}(1)}_k$ \\
\hline
Quantum structure & -- & Jacobian ring & $U_q(\widehat{\mathfrak{u}(1)})$ \\
\hline
\end{tabular}
\caption{Comparison of the three illustrative examples.}
\label{tab:examples_summary}
\end{table}
